% ============================================================
% MATHEMATICAL FOUNDATIONS OF MODERN AI
% Lecture Notes (5-day workshop) — English version
% Compilation: xelatex -shell-escape notes.tex (2 passes)
% ============================================================
\documentclass[12pt,a4paper,twoside]{book}

% --- Encoding and language (XeLaTeX) ---
\usepackage{fontspec}
\usepackage{polyglossia}
\setmainlanguage{english}

% --- Page layout ---
\usepackage[margin=2.5cm]{geometry}
\usepackage{fancyhdr}
\pagestyle{fancy}
\fancyhead[LE]{\leftmark}
\fancyhead[RO]{\rightmark}
\fancyfoot[C]{\thepage}
\setlength{\headheight}{14.5pt}

% --- Mathematics ---
\usepackage{amsmath,amssymb,amsthm,mathtools}
\usepackage{bm}

% --- Graphics ---
\usepackage{tikz,pgfplots}
\pgfplotsset{compat=1.18}
\usetikzlibrary{patterns,arrows.meta,calc,positioning,shapes,fit,backgrounds,chains}

% --- Code ---
\usepackage{minted}
\setminted{fontsize=\small,frame=lines,framesep=2mm,breaklines=true}
\setminted[python]{python3=true}

% --- Tables ---
\usepackage{booktabs,array,multirow,longtable}

% --- Colors and boxes ---
\usepackage[dvipsnames,svgnames]{xcolor}
\definecolor{aiblue}{RGB}{59,130,246}
\definecolor{warnorange}{RGB}{255,140,0}
\definecolor{deepgreen}{RGB}{22,163,74}
\usepackage{tcolorbox}
\tcbuselibrary{theorems,breakable,skins}

% --- Custom environments ---
\newtcolorbox{aitip}{
  colback=aiblue!5,colframe=aiblue!70,
  title={\faRobot\ Note},
  fonttitle=\bfseries,breakable
}
\newtcolorbox{warningbox}{
  colback=warnorange!5,colframe=warnorange!70,
  title={\faExclamationTriangle\ Caution},
  fonttitle=\bfseries,breakable
}
\newtcolorbox{exercise}{
  colback=gray!5,colframe=gray!50,
  title={Exercise},
  fonttitle=\bfseries,breakable
}
\newtcolorbox{intuition}{
  colback=deepgreen!5,colframe=deepgreen!70,
  title={\faLightbulb\ Intuition},
  fonttitle=\bfseries,breakable
}

% --- Links ---
\usepackage{hyperref}
\hypersetup{colorlinks=true,linkcolor=NavyBlue,citecolor=ForestGreen,urlcolor=Maroon}

% --- Icons ---
\usepackage{fontawesome5}

% --- Theorem environments ---
\newtheorem{definition}{Definition}[chapter]
\newtheorem{example}{Example}[chapter]
\newtheorem{proposition}{Proposition}[chapter]

% ============================================================
\title{
  \vspace{-2cm}
  {\Huge\bfseries Mathematical Foundations}\\[6pt]
  {\Huge\bfseries of Modern AI}\\[12pt]
  {\LARGE A 5-day workshop}\\[24pt]
  {\normalsize Yaé Ulrich Gaba}
}
\date{2026}
\author{}

% ============================================================
\setlength{\parindent}{0pt}
\setlength{\parskip}{0.5\baselineskip plus 0.1\baselineskip minus 0.1\baselineskip}

\begin{document}

\frontmatter
\maketitle

\tableofcontents
\newpage

% --- Preface ---
\chapter*{Preface}
\addcontentsline{toc}{chapter}{Preface}

Modern AI moves fast, but the ideas that actually power today's systems are stable and accessible. The same handful of mathematical structures --- vector spaces, gradients, probability distributions, inductive bias, generalization --- show up in every model we deploy, from a logistic regression to a diffusion model.

This 5-day workshop builds that mathematical intuition end-to-end. We start from refreshers in linear algebra and calculus, build toward the structures behind multilayer perceptrons, autoencoders, and diffusion models, and finish with the discipline of evaluating what we have built. The goal is not to turn participants into deep-learning specialists. It is to give them the mental structures to evaluate a method, talk to a technical team, frame a business problem as a learning problem, and decide when --- and when not --- to deploy an AI tool.

\paragraph{Who this is for.} Engineers, product managers, analysts, and technical team leads who make decisions about AI. Practitioners who work alongside ML teams and want a real grasp of the foundations. Researchers from other disciplines who want to move from talking about AI to using it.

\paragraph{What you will be able to do after the workshop.}
\begin{itemize}
  \item Read a description of a neural network architecture and understand what each piece is doing mathematically.
  \item Reason about why a model is trained the way it is --- and when a training choice signals a problem.
  \item Translate a real-world problem into the structured language of a learning problem.
  \item Identify the inductive bias built into a model and judge whether it matches the data.
  \item Recognize the standard failure modes (overfitting, distribution shift, calibration loss) and the standard remedies.
  \item Decide when AI is the right tool, when a simpler method is enough, and when more data or a different framing is the real fix.
\end{itemize}

\paragraph{Prerequisites.} Comfortable with high-school and first-year undergraduate mathematics (vectors, functions, derivatives, basic probability). Some Python familiarity helps for the labs but is not required for the lectures.

\paragraph{Format.} Five days. Each day pairs a lecture with a hands-on Jupyter notebook lab. All code runs on Google Colab (free tier with GPU) or any local Python 3.10+ environment with PyTorch and NumPy.

\vfill
\noindent Yaé Ulrich Gaba\\
Cotonou, 2026

% ============================================================
\mainmatter

\chapter{Linear algebra and parametric models}
\label{ch:day1}

\begin{quote}
\emph{``A neural network is a stack of linear maps with non-linear interruptions. The mathematics has been understood for a century; the engineering surprise is that it works at scale.''}
\end{quote}

% ============================================================
A modern AI model is, at its core, a function. It takes an input vector $\bm{x} \in \mathbb{R}^d$ --- pixels of an image, embeddings of a sentence, features of a customer --- and returns an output: a class label, a probability, another vector. The vocabulary of that function is linear algebra. Vector spaces describe where the data lives. Linear maps move data between spaces. The non-linear ingredients (ReLU, softmax) are punctuation. Almost everything else about a network --- depth, width, embedding dimension, latent space, attention heads --- is a choice of how to compose linear maps.

This chapter sets up the linear-algebraic spine of the workshop. We will not prove the spectral theorem. We will name the structures, give the geometric intuition, and write the small amount of code needed to see them in action. Day 2 will introduce gradients to make these models trainable; Day 3 will add probability to make them generative. But the structures are all here.

\section{Vector spaces, the short version}

A \emph{vector space} over $\mathbb{R}$ is a set $V$ where you can add elements and multiply them by real numbers, and the usual rules hold (associativity, distributivity, a zero element). Every concrete example we care about in this workshop is some flavor of $\mathbb{R}^d$: ordered tuples of real numbers.

\begin{definition}[The standard example]
$\mathbb{R}^d$ is the set of $d$-tuples $\bm{x} = (x_1, \ldots, x_d)$ with addition and scalar multiplication done component-wise.
\end{definition}

\begin{example}
A grayscale image of $28 \times 28$ pixels is an element of $\mathbb{R}^{784}$ once we flatten it. A sentence embedding from a small language model is an element of $\mathbb{R}^{768}$. A user's row in a recommender system's feature table is an element of $\mathbb{R}^{50}$ or so. These all look very different in the world but identical to the mathematics --- they are points in a vector space, and the model treats them as such.
\end{example}

The \emph{dimension} of a vector space is the number of independent directions you can move in. $\mathbb{R}^{784}$ has 784 of them, one per pixel. In practice, the data does not fill the space: nearly all $\mathbb{R}^{784}$ vectors look like random noise, and almost none of them look like a digit. This observation --- that real data lives on a much lower-dimensional subset --- is the \emph{manifold hypothesis}, and it is the reason embeddings work.

\begin{intuition}
A high-dimensional vector space is mostly empty. Modern AI exploits this: even when the ambient space has $10^4$ dimensions, the structure we care about (the set of plausible images, the set of meaningful sentences) lives on a tiny, curved subset. Learning that subset is a large part of what training does.
\end{intuition}

\section{Linear maps and matrices}

A \emph{linear map} $f: \mathbb{R}^d \to \mathbb{R}^k$ is a function that respects addition and scalar multiplication:
\[
f(\bm{x} + \bm{y}) = f(\bm{x}) + f(\bm{y}), \qquad f(\alpha \bm{x}) = \alpha f(\bm{x}).
\]
Every such map is given by a matrix $W \in \mathbb{R}^{k \times d}$: $f(\bm{x}) = W \bm{x}$. Pick a basis and the map becomes a list of numbers.

\begin{definition}[Affine map]
An affine map is a linear map plus a constant offset: $\bm{x} \mapsto W \bm{x} + \bm{b}$, where $\bm{b} \in \mathbb{R}^k$ is the bias. Every layer of a standard neural network is an affine map followed by a non-linearity.
\end{definition}

The matrix $W$ does three things at once: it scales, it rotates, and it projects onto a (possibly lower-dimensional) subspace. The singular value decomposition makes this precise --- any $W$ can be written $U \Sigma V^\top$ where $U$ and $V$ are rotations and $\Sigma$ is a non-negative diagonal scaling. We will not need the SVD machinery in detail, but the picture is worth holding:

\begin{intuition}
Every linear layer in a neural network rotates the input space, stretches some directions and shrinks others, and possibly throws some directions away (those with singular value zero). Stacking layers stacks these operations. The non-linearities are what prevent the stack from collapsing into a single linear map.
\end{intuition}

\section{From the perceptron to the MLP}

The simplest neural network is a single neuron, the \emph{perceptron} (Rosenblatt, 1958):
\[
y = \sigma(\bm{w}^\top \bm{x} + b)
\]
where $\bm{w} \in \mathbb{R}^d$ is a vector of weights, $b \in \mathbb{R}$ is a bias, and $\sigma$ is a non-linear function. With $\sigma(z) = 1/(1 + e^{-z})$ (the logistic sigmoid), this is exactly logistic regression dressed up in different vocabulary.

\begin{proposition}[Why a single layer is not enough]
A single affine layer can only separate input space by a hyperplane. The XOR function --- output is 1 if exactly one of two binary inputs is 1 --- cannot be expressed by any single perceptron. This is the original 1969 critique of Minsky and Papert that briefly killed the field.
\end{proposition}

The fix is composition. A \emph{multilayer perceptron} (MLP) stacks affine maps with non-linearities between them. A two-layer MLP from $\mathbb{R}^d$ to $\mathbb{R}^k$ is:
\[
\bm{h} = \sigma(W_1 \bm{x} + \bm{b}_1), \qquad \bm{y} = W_2 \bm{h} + \bm{b}_2
\]
where $W_1 \in \mathbb{R}^{m \times d}$, $W_2 \in \mathbb{R}^{k \times m}$, and $m$ is the width of the hidden layer. The non-linearity $\sigma$ is applied component-wise. Modern networks use $\sigma(z) = \max(0, z)$, the rectified linear unit (ReLU), almost everywhere.

\begin{aitip}
The universal approximation theorem (Cybenko, 1989; Hornik, 1991) says that even a one-hidden-layer MLP with enough width can approximate any continuous function on a compact set to arbitrary accuracy. This is a statement about \emph{existence}, not about \emph{learnability} or efficiency. Modern deep networks are not deep because shallow networks cannot represent the function. They are deep because depth makes the right functions easier to find by gradient descent.
\end{aitip}

\subsection{Counting parameters}

Take an MLP with input dimension $d = 784$, one hidden layer of width $m = 256$, and output dimension $k = 10$. The parameter count is:
\[
\underbrace{784 \cdot 256 + 256}_{\text{first layer: } W_1, \bm{b}_1} + \underbrace{256 \cdot 10 + 10}_{\text{second layer: } W_2, \bm{b}_2} = 203\,530.
\]
A small MNIST classifier already has a couple hundred thousand parameters. GPT-3 has 175 billion. The numbers grow fast, but the structure does not change --- every layer is still an affine map followed by a non-linearity.

\subsection{The MLP in code}

\begin{minted}{python}
import torch
import torch.nn as nn

class TwoLayerMLP(nn.Module):
    def __init__(self, d_in: int, d_hidden: int, d_out: int):
        super().__init__()
        self.fc1 = nn.Linear(d_in, d_hidden)   # affine map: W1 x + b1
        self.fc2 = nn.Linear(d_hidden, d_out)  # affine map: W2 h + b2

    def forward(self, x: torch.Tensor) -> torch.Tensor:
        h = torch.relu(self.fc1(x))            # non-linearity
        y = self.fc2(h)
        return y

model = TwoLayerMLP(d_in=784, d_hidden=256, d_out=10)
n_params = sum(p.numel() for p in model.parameters())
print(f"Number of parameters: {n_params:,}")
\end{minted}

The \texttt{nn.Linear} layer holds exactly the $W$ and $\bm{b}$ from the formulas above. Everything else in this workshop is variations on this theme: more layers, different non-linearities, weight-sharing (CNNs), structured connections (GNNs), attention (Transformers). The base unit is the affine map.

\begin{exercise}
A residual block in a ResNet has the form $\bm{y} = \bm{x} + F(\bm{x})$, where $F$ is a small sub-network. Argue informally why this should make training easier than the bare composition $\bm{y} = F(\bm{x})$. (Hint: what happens if $F$ outputs zero? What does that imply for the gradient?)
\end{exercise}

\section{Embeddings: vectors as representations}

Up to this point, we have treated the input $\bm{x}$ as if it were already a vector. But the real input is often something else: a word, a user ID, a discrete category. Turning it into a vector is the first decision the model makes.

\begin{definition}[Embedding]
An embedding is a function from a discrete set $V$ (vocabulary, set of users, set of categories) into a vector space $\mathbb{R}^d$. Concretely it is a matrix $E \in \mathbb{R}^{|V| \times d}$: each row is the vector representation of one item.
\end{definition}

The embedding matrix is a parameter of the model. It is learned. After training, similar items end up close in the embedding space, and the geometry of that space carries the meaning the model has discovered.

\begin{example}
Word2Vec (Mikolov et al., 2013) trains an embedding matrix for English words so that the cosine similarity of two word vectors tracks their semantic similarity. The famous geometric regularity $\bm{e}_{\text{king}} - \bm{e}_{\text{man}} + \bm{e}_{\text{woman}} \approx \bm{e}_{\text{queen}}$ is an empirical observation about the resulting space, not an objective the model was trained on. The training objective was much simpler: predict a word from its context.
\end{example}

In modern language models, the input embedding is the first layer (a lookup in an embedding matrix), and the rest of the network operates on those vectors. The trick is that the same vector space is doing double duty: it is the input space \emph{and} the model's working memory.

\begin{warningbox}
``Embedding'' is overloaded in machine learning. It can mean a learned input representation (Word2Vec, an LLM's token embedding), a hidden activation inside a network (``the embedding of this image is the output of the second-to-last layer''), or a separately trained representation used as a feature in a downstream model (sentence-transformers). All three are vectors. The differences are about how they are produced and what they are used for.
\end{warningbox}

\section{Autoencoders and the geometry of latent space}

An \emph{autoencoder} learns a compressed representation of its input by trying to reproduce that input from the compression. It has two parts:
\[
\underbrace{\bm{z} = f_{\text{enc}}(\bm{x})}_{\text{encoder, } \mathbb{R}^d \to \mathbb{R}^k} \qquad \underbrace{\hat{\bm{x}} = f_{\text{dec}}(\bm{z})}_{\text{decoder, } \mathbb{R}^k \to \mathbb{R}^d}
\]
with $k < d$. The training objective is reconstruction: minimize $\|\bm{x} - \hat{\bm{x}}\|^2$ over a dataset. Because the bottleneck $\bm{z}$ is lower-dimensional than $\bm{x}$, the autoencoder cannot memorize the input; it has to find structure.

\begin{intuition}
PCA is a linear autoencoder. If $f_{\text{enc}}$ and $f_{\text{dec}}$ are forced to be linear maps and we minimize squared reconstruction error, the optimal solution projects onto the top-$k$ principal components. Non-linear autoencoders generalize this idea: they learn a curved $k$-dimensional surface inside $\mathbb{R}^d$ that fits the data.
\end{intuition}

The bottleneck space $\mathbb{R}^k$ is called the \emph{latent space}. It is where the model's compressed understanding of the data lives. Two questions are usually worth asking about any latent space:

\begin{itemize}
  \item \textbf{What does each dimension mean?} Often nothing on its own --- the axes are arbitrary --- but directions in the space often correspond to semantic factors (lighting, pose, identity for a face dataset; topic, sentiment, tense for a text dataset).
  \item \textbf{How does it transport between nearby points?} If small movements in $\bm{z}$ produce small, plausible changes in $\hat{\bm{x}}$, the model has learned the right manifold. If small movements produce jumps or nonsense, it has not.
\end{itemize}

Day 3 will return to autoencoders in their probabilistic form (variational autoencoders, VAEs) as one of the bridges to generative modelling. The autoencoder structure --- encoder, latent, decoder --- is also exactly the structure of a diffusion model in disguise.

\section{What you have seen}

Every parametric model in this workshop is built from three pieces: vectors (the data and intermediate representations), affine maps (the layers, with their weights and biases), and non-linearities (the activations that prevent collapse). Modern architectures --- CNNs, RNNs, Transformers, GNNs, diffusion models --- are variations on which affine maps are tied together and which non-linearities are inserted where. The names are different. The mathematics is the same.

What we have not yet done is make any of this \emph{learn}. So far the weights $W$ and biases $\bm{b}$ are just numbers; the model is a function. Day 2 introduces the gradient, the loss function, and the optimization machinery that turns this function into something that can be fit to data. The structure stays the same. Training is what fills it in.

\section*{Lab}

Open the Day 1 notebook (\texttt{Day1\_LinearAlgebra\_ParametricModels.ipynb}). You will:

\begin{enumerate}
  \item Implement a two-layer MLP by hand in NumPy --- writing the matrix multiplications and ReLU explicitly --- and verify the parameter count matches what we computed above.
  \item Train the same MLP in PyTorch on MNIST, getting roughly 97\% test accuracy in under a minute.
  \item Extract the activations of the hidden layer for the test set, project them to 2D with PCA, and color the points by their true digit class. You will see clusters --- the network has learned a representation in which similar digits live near each other.
  \item Train a small autoencoder on MNIST with a 2-dimensional latent space. Visualize the latent space directly and walk along straight lines in it: you will see digits morph smoothly into other digits.
\end{enumerate}

By the end of the lab you will have built the machinery that the rest of the workshop assumes: a working MLP, a working autoencoder, and the habit of inspecting what a network has learned by looking at its internal representations.

\chapter{Calculus and optimization}
\label{ch:day2}

\begin{quote}
\emph{``Training a neural network is the chain rule, applied carefully, many times per second.''}
\end{quote}

% ============================================================
At the end of Day 1 we had a model: an MLP, a stack of affine maps and ReLUs, with weights $W$ and biases $\bm{b}$ that were just random numbers. The model was a function. It did nothing useful. To make it useful we need to adjust those numbers so the model's outputs match the data --- and we need to do that adjustment automatically, because there are too many numbers to set by hand.

This chapter is about that adjustment. The mathematics is calculus: derivatives say how a function's output changes when you nudge its inputs, and gradients generalize that idea to functions of many variables. The algorithm is gradient descent: take a small step in the direction that decreases the loss. The engineering trick that makes it tractable for million-parameter networks is the chain rule, applied recursively --- backpropagation. By the end of the day you will have implemented backpropagation by hand, in NumPy, for the MLP we built yesterday.

\section{The loss function}

Before we can optimize anything we have to say what we are optimizing. We need a single real number that measures \emph{how wrong} the model is on a dataset. That number is the \emph{loss}.

\begin{definition}[Loss function]
Given a model $f_\theta$ with parameters $\theta$ and a dataset $\{(\bm{x}_i, y_i)\}_{i=1}^N$, a loss function $\mathcal{L}(\theta)$ is a real-valued function of $\theta$ that is small when the model fits the data well and large when it does not. Training is the optimization problem $\min_\theta \mathcal{L}(\theta)$.
\end{definition}

Two losses cover most of this workshop:

\paragraph{Mean squared error (regression).} For real-valued targets:
\[
\mathcal{L}_{\text{MSE}}(\theta) = \frac{1}{N} \sum_{i=1}^N \|f_\theta(\bm{x}_i) - y_i\|^2.
\]
This is the squared distance between prediction and target, averaged over the dataset.

\paragraph{Cross-entropy (classification).} For class labels $y_i \in \{1, \ldots, K\}$, the model outputs a probability distribution $f_\theta(\bm{x}_i) \in \Delta^{K-1}$ over classes (via softmax), and:
\[
\mathcal{L}_{\text{CE}}(\theta) = -\frac{1}{N} \sum_{i=1}^N \log f_\theta(\bm{x}_i)_{y_i}.
\]
This is the negative log-likelihood under the model's predicted distribution. Day 3 explains why \emph{this} loss and not some other.

\begin{intuition}
The loss is a landscape. The parameters $\theta$ are coordinates. Training is the act of finding a low point. The landscape has millions of dimensions (one per parameter) but is otherwise just a function $\mathcal{L}: \mathbb{R}^P \to \mathbb{R}$. Everything in this chapter is about how to walk downhill on that landscape efficiently.
\end{intuition}

\section{Derivatives and gradients}

For a scalar function $f: \mathbb{R} \to \mathbb{R}$, the derivative $f'(x)$ is the rate of change at $x$: if we nudge $x$ by $\epsilon$, the output changes by approximately $f'(x) \cdot \epsilon$. The gradient generalizes this to functions of many variables.

\begin{definition}[Gradient]
For a function $f: \mathbb{R}^d \to \mathbb{R}$, the gradient at $\bm{x}$ is the vector of partial derivatives:
\[
\nabla f(\bm{x}) = \left( \frac{\partial f}{\partial x_1}(\bm{x}), \, \frac{\partial f}{\partial x_2}(\bm{x}), \, \ldots, \, \frac{\partial f}{\partial x_d}(\bm{x}) \right).
\]
The gradient points in the direction of steepest \emph{ascent}. The opposite direction, $-\nabla f(\bm{x})$, is steepest descent.
\end{definition}

This is the key geometric fact behind every optimization algorithm in this workshop: to decrease a function, follow its negative gradient.

\begin{example}
Let $f(x_1, x_2) = x_1^2 + 3 x_2^2$. The gradient is $\nabla f = (2 x_1, 6 x_2)$. At the point $(1, 1)$ the gradient is $(2, 6)$. The function increases most rapidly in the direction $(2, 6)/\|(2, 6)\|$ and decreases most rapidly in the opposite direction. Both axes point outward from the minimum at $(0, 0)$, but the $x_2$ direction is steeper, because the function curves more steeply along $x_2$.
\end{example}

\section{Gradient descent}

The simplest optimization algorithm is the one suggested by the previous section: start somewhere, compute the gradient, step in the negative-gradient direction, repeat.

\begin{definition}[Gradient descent]
Given a learning rate $\eta > 0$ and an initial point $\theta_0$, gradient descent produces a sequence
\[
\theta_{t+1} = \theta_t - \eta \, \nabla \mathcal{L}(\theta_t).
\]
\end{definition}

The learning rate $\eta$ controls how big a step we take. If $\eta$ is too small, training is slow. If $\eta$ is too large, the steps overshoot and the loss oscillates or diverges.

\begin{warningbox}
The learning rate is the single most consequential hyperparameter in deep learning. ``My model is not training'' is, in practice, almost always ``my learning rate is wrong.'' The lab today includes a learning-rate sweep so you can see this directly.
\end{warningbox}

\section{The chain rule and backpropagation}

For an MLP with two layers, the loss is a composition of functions:
\[
\mathcal{L}(\theta) = \ell\bigl(f_2(f_1(\bm{x}))\bigr)
\]
where $f_1$ is the first affine-plus-ReLU layer, $f_2$ is the second, and $\ell$ is the loss applied to the output. To do gradient descent on $\mathcal{L}$ we need the gradient with respect to every parameter in $f_1$ and $f_2$.

The chain rule from single-variable calculus generalizes to vector-valued functions. If $\bm{y} = g(\bm{u})$ and $\bm{u} = h(\bm{x})$, then
\[
\frac{\partial \bm{y}}{\partial \bm{x}} = \frac{\partial \bm{y}}{\partial \bm{u}} \cdot \frac{\partial \bm{u}}{\partial \bm{x}}
\]
where the partial derivatives are Jacobian matrices and the product is matrix multiplication. Applied repeatedly to the composition above, this lets us compute $\partial \mathcal{L} / \partial \theta$ for every parameter.

\begin{intuition}
The forward pass computes the loss from inputs to scalar output. The backward pass computes derivatives in the opposite direction: it starts from the scalar loss (whose derivative with respect to itself is 1) and propagates partial derivatives back through every layer, multiplying Jacobians as it goes. The total work of the backward pass is roughly the same as the forward pass --- a factor of 2 to 3, depending on the network. This is why training is feasible at all.
\end{intuition}

\subsection{Backprop for a 2-layer MLP, by hand}

For the two-layer MLP from Day 1:
\begin{align*}
\bm{h}_{\text{pre}} &= W_1 \bm{x} + \bm{b}_1 \\
\bm{h} &= \mathrm{ReLU}(\bm{h}_{\text{pre}}) \\
\bm{y} &= W_2 \bm{h} + \bm{b}_2
\end{align*}
with cross-entropy loss after a softmax, the gradients have a closed form. Let $\bm{p} = \mathrm{softmax}(\bm{y})$ and $\bm{y}^*$ be the one-hot target. Then:
\begin{align*}
\frac{\partial \mathcal{L}}{\partial \bm{y}} &= \bm{p} - \bm{y}^* &\text{(softmax-CE collapses)} \\
\frac{\partial \mathcal{L}}{\partial W_2} &= \frac{\partial \mathcal{L}}{\partial \bm{y}} \cdot \bm{h}^\top, \quad \frac{\partial \mathcal{L}}{\partial \bm{b}_2} = \frac{\partial \mathcal{L}}{\partial \bm{y}} \\
\frac{\partial \mathcal{L}}{\partial \bm{h}} &= W_2^\top \cdot \frac{\partial \mathcal{L}}{\partial \bm{y}} \\
\frac{\partial \mathcal{L}}{\partial \bm{h}_{\text{pre}}} &= \frac{\partial \mathcal{L}}{\partial \bm{h}} \odot \mathbf{1}[\bm{h}_{\text{pre}} > 0] &\text{(ReLU gradient)} \\
\frac{\partial \mathcal{L}}{\partial W_1} &= \frac{\partial \mathcal{L}}{\partial \bm{h}_{\text{pre}}} \cdot \bm{x}^\top, \quad \frac{\partial \mathcal{L}}{\partial \bm{b}_1} = \frac{\partial \mathcal{L}}{\partial \bm{h}_{\text{pre}}}.
\end{align*}
Five lines. That is backpropagation for a two-layer MLP. The lab implements these formulas in NumPy and checks them against PyTorch's autograd to the bit.

\subsection{Why we never write this code in practice}

For a two-layer MLP we could derive these formulas. For a fifty-layer ResNet, or for a Transformer with cross-attention, the algebra is hopeless. Modern frameworks (PyTorch, JAX, TensorFlow) record the computational graph of the forward pass and replay it backward, applying the chain rule at each node. Every primitive operation (matmul, ReLU, softmax) has a hand-coded backward; the framework composes them automatically. This is what \emph{differentiable programming} means.

\begin{aitip}
You will never write the backward pass for a production model. But you should understand what it is, because almost every training failure --- vanishing gradients, exploding gradients, dead ReLUs, gradient checkpointing trade-offs --- is a story about backpropagation. ``My loss is NaN'' is a story about backpropagation.
\end{aitip}

\section{Stochastic gradient descent}

The loss is a sum over the dataset:
\[
\mathcal{L}(\theta) = \frac{1}{N} \sum_{i=1}^N \ell_i(\theta).
\]
Computing $\nabla \mathcal{L}$ exactly requires a full pass over the dataset. For MNIST that is 60\,000 examples per step. For ImageNet, 1.2 million. For an LLM pretrained on $10^{12}$ tokens, the number is meaningless.

\begin{definition}[Stochastic gradient descent, SGD]
At each step, sample a small batch $B \subset \{1, \ldots, N\}$ and update with the batch gradient:
\[
\theta_{t+1} = \theta_t - \eta \cdot \frac{1}{|B|} \sum_{i \in B} \nabla \ell_i(\theta_t).
\]
The batch gradient is a noisy but unbiased estimate of the full gradient.
\end{definition}

Two things make SGD work in practice:

\begin{enumerate}
  \item \textbf{Computational.} A batch of 128 fits in memory, runs in milliseconds, and gives a usable step. Full-batch gradient descent would be a hundred times slower per step for the same compute budget.
  \item \textbf{Statistical.} The noise in the gradient estimate acts as implicit regularization, helping the optimizer escape narrow minima. This is part of why neural networks generalize at all (Hardt et al., 2016; Keskar et al., 2017).
\end{enumerate}

\section{Adam: momentum and adaptive step sizes}

Plain SGD is the simplest optimizer, but rarely the best one in practice. Adam (Kingma \& Ba, 2015) adds two ingredients:

\paragraph{Momentum.} Instead of taking a step in the direction of the current gradient, take a step in the direction of an exponentially-weighted average of recent gradients. This smooths out noise and accelerates progress along consistent directions.

\paragraph{Per-parameter learning rates.} Each parameter gets its own effective learning rate, based on a running estimate of the magnitude of \emph{its} recent gradients. Parameters with consistently small gradients take larger steps; parameters with consistently large gradients take smaller steps.

The update rule, with $\beta_1 = 0.9$, $\beta_2 = 0.999$, $\epsilon = 10^{-8}$:
\begin{align*}
\bm{m}_t &= \beta_1 \bm{m}_{t-1} + (1 - \beta_1) \nabla \mathcal{L}(\theta_t) &\text{(momentum)} \\
\bm{v}_t &= \beta_2 \bm{v}_{t-1} + (1 - \beta_2) (\nabla \mathcal{L}(\theta_t))^2 &\text{(squared-gradient running mean)} \\
\theta_{t+1} &= \theta_t - \eta \cdot \frac{\hat{\bm{m}}_t}{\sqrt{\hat{\bm{v}}_t} + \epsilon}
\end{align*}
where $\hat{\bm{m}}, \hat{\bm{v}}$ are bias-corrected versions of $\bm{m}, \bm{v}$. The square root and division are element-wise.

Adam is the default optimizer for almost every modern network. SGD with momentum still beats it for some computer-vision benchmarks (Wilson et al., 2017), but Adam's tolerance of bad hyperparameter choices is what makes it the practical default.

\section{Reading a learning curve}

A learning curve plots loss (or accuracy) against training step. It is the diagnostic instrument of deep learning. Four characteristic shapes:

\begin{itemize}
  \item \textbf{Healthy:} train loss decreases smoothly, validation loss decreases too and tracks within a small gap. The model is learning, and what it learns generalizes.
  \item \textbf{Overfitting:} train loss decreases, validation loss decreases for a while and then starts to climb. The model is memorizing training examples instead of learning the underlying pattern. Day 5 returns to this.
  \item \textbf{Underfitting:} both losses plateau high. The model lacks capacity, the optimizer is stuck, or the learning rate is wrong.
  \item \textbf{Diverging:} loss spikes upward, often to infinity. The learning rate is too large, or the loss is being computed on a non-differentiable region (log of zero, division by zero in attention).
\end{itemize}

\begin{intuition}
Reading the learning curve is the first thing to do when training does not work. Most pathologies have a characteristic shape, and the curve tells you which one you are looking at within seconds. Today's lab generates each of these shapes deliberately so you can recognize them next time.
\end{intuition}

\section{What you have seen}

Training a neural network is gradient descent on a high-dimensional loss landscape, made tractable by the chain rule (backpropagation) and made practical by stochastic mini-batches and adaptive optimizers like Adam. The mathematics is calculus; the engineering is automatic differentiation. Every modern framework gives you the second for free, but only the first lets you reason about what is going wrong when something does.

Day 3 adds the missing ingredient: probability. We have been speaking of ``predicting'' outputs and ``matching'' targets, but generative models do something stranger --- they learn an entire probability distribution and sample from it. That requires us to take the loss function seriously as a likelihood.

\section*{Lab}

Open the Day 2 notebook (\texttt{Day2\_Calculus\_Optimization.ipynb}). You will:

\begin{enumerate}
  \item Compute the gradient of a small function by hand, then verify against PyTorch's autograd.
  \item Add a manual backward pass to the NumPy MLP from Day 1, then check every gradient against PyTorch's autograd to the last decimal --- the standard gradient-check trick used to debug new layers.
  \item Train the manual MLP on MNIST with vanilla SGD and your own training loop. No frameworks.
  \item Replace the manual training loop with PyTorch + Adam and observe the difference.
  \item Sweep the learning rate over $\{10^{-1}, 10^{-2}, 10^{-3}, 10^{-4}\}$ and plot the resulting learning curves on a single axis. You will see all four characteristic shapes from above.
\end{enumerate}

By the end of the lab you will have written and verified a backward pass, trained a network from first principles, and built the habit of looking at the learning curve \emph{before} blaming the model.

\chapter{Probability and generative modelling}
\label{ch:day3}

\begin{quote}
\emph{``A generative model is a probability distribution you can sample from. Everything else --- the loss function, the architecture, the training algorithm --- is in service of learning that distribution from data.''}
\end{quote}

% ============================================================
Days 1 and 2 built a function: an MLP that maps inputs to outputs, with weights trained by gradient descent on a loss. The setup was deterministic. A given input produced a given output. That picture is enough for classification and regression, but it does not capture what modern generative models do. A diffusion model does not predict a single output from an input; it samples one image (or many different images) from a learned distribution. A language model does not predict one next token; it gives you a distribution over the vocabulary, and the sampler draws from it.

To talk about that, we need probability. This chapter is the bridge: it reintroduces the loss function from Day 2 as a likelihood, develops the variational autoencoder as the simplest non-trivial generative model, and ends with the denoising diffusion picture that underlies essentially every state-of-the-art image and audio generator built in the last five years.

\section{Distributions, densities, and what we are trying to learn}

A probability distribution assigns weight to every possible outcome. For a finite set, a distribution is a list of non-negative numbers summing to 1. For a continuous space like $\mathbb{R}^d$, it is a density function $p(\bm{x}) \geq 0$ whose integral over the whole space is 1.

\begin{definition}[Density]
A probability density $p: \mathbb{R}^d \to \mathbb{R}_{\geq 0}$ satisfies $\int p(\bm{x}) \, d\bm{x} = 1$. For any region $A \subseteq \mathbb{R}^d$, the probability that a sample falls in $A$ is $\int_A p(\bm{x}) \, d\bm{x}$.
\end{definition}

The data we care about --- images, texts, molecules --- comes from some unknown distribution $p_{\text{data}}$. Generative modelling means: estimate $p_{\text{data}}$ from samples, well enough to generate new samples that look like they came from it.

\begin{intuition}
For MNIST, $p_{\text{data}}$ is a density on $\mathbb{R}^{784}$. Almost every point in $\mathbb{R}^{784}$ is a noise image with density essentially zero. Real digits live on a tiny, curved subset (the manifold from Day 1), where the density is enormous. A generative model is something that has internalized the shape of that high-density region.
\end{intuition}

\section{Likelihood: a loss with a meaning}

Suppose we have a parametric family of densities $\{p_\theta : \theta \in \Theta\}$. Given samples $\bm{x}_1, \ldots, \bm{x}_N$ drawn from $p_{\text{data}}$, the \emph{likelihood} is the probability density that the model assigns to the data:
\[
\mathcal{L}_{\text{lik}}(\theta) = \prod_{i=1}^N p_\theta(\bm{x}_i).
\]
\emph{Maximum likelihood estimation} (MLE) picks the $\theta$ that makes the data look most plausible under $p_\theta$. In practice we work with the negative log-likelihood, because it turns the product into a sum and gives us something to minimize:
\[
- \log \mathcal{L}_{\text{lik}}(\theta) = - \sum_{i=1}^N \log p_\theta(\bm{x}_i).
\]

\begin{proposition}[Cross-entropy is negative log-likelihood]
For a classifier with softmax output $p_\theta(y \mid \bm{x})$, the cross-entropy loss from Day 2,
\[
\mathcal{L}_{\text{CE}}(\theta) = - \frac{1}{N} \sum_{i=1}^N \log p_\theta(y_i \mid \bm{x}_i),
\]
is exactly $-\frac{1}{N} \log \mathcal{L}_{\text{lik}}(\theta)$. Training a classifier with cross-entropy is maximum likelihood on the conditional distribution $p_\theta(y \mid \bm{x})$.
\end{proposition}

The classifier from Day 2 is already a generative model --- of labels conditional on inputs. To generate \emph{inputs} we need a model of the input distribution itself. That is what the rest of this chapter is about.

\section{KL divergence: distance between distributions}

To compare two distributions we use the Kullback--Leibler divergence:
\[
D_{\mathrm{KL}}(q \,\|\, p) = \int q(\bm{x}) \log \frac{q(\bm{x})}{p(\bm{x})} \, d\bm{x} = \mathbb{E}_{\bm{x} \sim q}\!\left[ \log q(\bm{x}) - \log p(\bm{x}) \right].
\]
KL is non-negative, zero if and only if $p = q$, and asymmetric. It is not a metric --- $D_{\mathrm{KL}}(p \,\|\, q) \neq D_{\mathrm{KL}}(q \,\|\, p)$ in general --- but it is the natural object when one of the distributions is the data and the other is the model.

\begin{proposition}[MLE minimizes KL to the data]
Up to a constant that does not depend on $\theta$, the negative log-likelihood equals $D_{\mathrm{KL}}(p_{\text{data}} \,\|\, p_\theta)$. Maximum likelihood is approximate KL minimization between the data distribution and the model distribution.
\end{proposition}

This is the cleanest statement of what training a generative model is doing: making the model distribution close to the data distribution in the KL sense. The rest of the chapter is two stories about \emph{how} to do that minimization when $p_\theta$ is a neural network.

\section{Variational autoencoders}

Recall the autoencoder from Day 1: an encoder $f_{\text{enc}}: \mathbb{R}^d \to \mathbb{R}^k$ and a decoder $f_{\text{dec}}: \mathbb{R}^k \to \mathbb{R}^d$, with $k < d$, trained to minimize $\|\bm{x} - f_{\text{dec}}(f_{\text{enc}}(\bm{x}))\|^2$. The autoencoder learns a latent space, but it has no notion of probability --- you cannot sample from it. The variational autoencoder (VAE, Kingma \& Welling, 2014) is the probabilistic cousin.

\subsection{The model}

The VAE assumes data is generated as follows:
\begin{enumerate}
  \item Sample a latent vector $\bm{z}$ from a simple prior, typically $p(\bm{z}) = \mathcal{N}(\bm{0}, I)$.
  \item Pass $\bm{z}$ through a neural network decoder $f_{\text{dec}}$ to get parameters of a distribution $p_\theta(\bm{x} \mid \bm{z})$ over data (e.g., a Gaussian with mean $f_{\text{dec}}(\bm{z})$).
  \item Sample $\bm{x}$ from $p_\theta(\bm{x} \mid \bm{z})$.
\end{enumerate}
The model distribution is the marginal $p_\theta(\bm{x}) = \int p_\theta(\bm{x} \mid \bm{z}) \, p(\bm{z}) \, d\bm{z}$. This integral is intractable for a neural-network decoder, so we cannot evaluate $\log p_\theta(\bm{x})$ directly. The VAE introduces a second neural network --- an \emph{encoder} $q_\phi(\bm{z} \mid \bm{x})$, typically Gaussian with mean and variance output by $f_{\text{enc}}$ --- and optimizes a lower bound on the log-likelihood.

\subsection{The Evidence Lower BOund (ELBO)}

For any choice of $q_\phi$,
\[
\log p_\theta(\bm{x}) \geq \underbrace{\mathbb{E}_{q_\phi(\bm{z} \mid \bm{x})}\!\left[ \log p_\theta(\bm{x} \mid \bm{z}) \right]}_{\text{reconstruction term}} - \underbrace{D_{\mathrm{KL}}(q_\phi(\bm{z} \mid \bm{x}) \,\|\, p(\bm{z}))}_{\text{regularizer to the prior}}.
\]
This is the ELBO. Maximizing the ELBO with respect to $\theta$ and $\phi$ jointly is the training objective. The first term wants the decoder to reconstruct $\bm{x}$ well from a $\bm{z}$ drawn from the encoder. The second term wants the encoder's posterior to look like the prior.

\begin{intuition}
The VAE is an autoencoder with two changes. First, the encoder outputs a distribution, not a single $\bm{z}$ --- you sample to get the latent. Second, an extra loss term pulls the encoder distributions toward the standard Gaussian prior. Together these changes make the latent space \emph{generative}: samples from the prior, decoded, look like data, because the encoder has been pushed to use the same region of latent space the prior covers.
\end{intuition}

\subsection{The reparameterization trick}

The expectation $\mathbb{E}_{q_\phi(\bm{z} \mid \bm{x})}[\cdot]$ is taken over a random $\bm{z}$, and the encoder produces the parameters of that randomness. We cannot directly differentiate through random sampling. The trick (Kingma \& Welling, 2014): write $\bm{z} = \bm{\mu}_\phi(\bm{x}) + \bm{\sigma}_\phi(\bm{x}) \odot \bm{\epsilon}$ with $\bm{\epsilon} \sim \mathcal{N}(\bm{0}, I)$. Now the randomness is in $\bm{\epsilon}$, which has no parameters; $\bm{\mu}_\phi$ and $\bm{\sigma}_\phi$ are deterministic and differentiable. Backprop flows through.

The lab today trains a small VAE on MNIST and visualizes both the latent space and what happens when you decode a smooth path through it.

\section{Denoising diffusion}

Diffusion models (Sohl-Dickstein et al., 2015; Ho et al., 2020) take a different route to the same destination. Instead of trying to learn the data distribution directly, they learn to \emph{reverse} a fixed noising process.

\subsection{The forward process}

Start from a real image $\bm{x}_0$. Define a sequence of corrupted versions $\bm{x}_1, \bm{x}_2, \ldots, \bm{x}_T$ by adding a small amount of Gaussian noise at each step:
\[
\bm{x}_t = \sqrt{1 - \beta_t} \, \bm{x}_{t-1} + \sqrt{\beta_t} \, \bm{\epsilon}_t, \qquad \bm{\epsilon}_t \sim \mathcal{N}(\bm{0}, I).
\]
The variances $\beta_t$ are small and chosen so that by $t = T$ (often $T = 1000$), $\bm{x}_T$ is essentially pure Gaussian noise. This forward process has no parameters; it is just a noise schedule.

A useful closed-form consequence: for any $t$, we can write $\bm{x}_t$ directly in terms of $\bm{x}_0$ and a single noise vector:
\[
\bm{x}_t = \sqrt{\bar{\alpha}_t} \, \bm{x}_0 + \sqrt{1 - \bar{\alpha}_t} \, \bm{\epsilon}, \quad \bm{\epsilon} \sim \mathcal{N}(\bm{0}, I),
\]
where $\bar{\alpha}_t = \prod_{s=1}^t (1 - \beta_s)$. This is what makes training efficient --- no need to simulate the chain step by step.

\subsection{The reverse process}

We train a neural network $\bm{\epsilon}_\theta(\bm{x}_t, t)$ that takes a noisy image and a timestep and predicts the noise that was added. The training loss is plain MSE:
\[
\mathcal{L}_{\text{diff}}(\theta) = \mathbb{E}_{\bm{x}_0, t, \bm{\epsilon}} \!\left[ \|\bm{\epsilon} - \bm{\epsilon}_\theta(\bm{x}_t, t)\|^2 \right]
\]
where $\bm{x}_t$ is the noisy version of $\bm{x}_0$ at timestep $t$ using the closed-form formula above. Once the network can predict the noise, we can subtract it: given a noisy $\bm{x}_t$, recover an estimate of $\bm{x}_0$ (or, more carefully, of $\bm{x}_{t-1}$).

\begin{aitip}
The training objective is just MSE on noise prediction. There is no log-likelihood, no KL divergence in the loss, no variational bound to negotiate. Diffusion's empirical advantage over VAEs is largely that this loss is much better behaved than the ELBO. The deep reason it works comes from the connection to score matching (Hyvärinen, 2005; Song \& Ermon, 2019): the noise prediction is, up to a known scaling, an estimate of the gradient of the log-density of the data distribution at the noised point.
\end{aitip}

\subsection{Sampling}

To generate a new image, start from pure noise $\bm{x}_T \sim \mathcal{N}(\bm{0}, I)$ and run the reverse process: at each step, use the network to predict the noise, then take one step closer to a clean image. After $T$ steps you have a sample from the model's approximation to $p_{\text{data}}$. Modern samplers (DDIM, Karras et al.) reduce the number of required steps from thousands to tens, at little loss of quality.

\begin{intuition}
A trained diffusion model is a denoiser that knows what data looks like at every noise level. Sampling is just denoising in reverse, starting from pure noise and ending at something the model considers a valid sample. The lab today trains this picture from scratch on a 2D spiral so you can watch the noise become a spiral in real time.
\end{intuition}

\section{Why probability changes the picture}

The deterministic networks of Days 1 and 2 were function approximators. The probabilistic networks of Day 3 are distribution approximators. The difference is consequential:

\begin{itemize}
  \item A classifier tells you the probability of each label. A diffusion model tells you the probability of each possible image.
  \item A classifier is queried once per input. A generative model is sampled --- it produces something different each time you call it.
  \item A classifier's failure is wrong predictions. A generative model's failure is producing samples that the data distribution would never produce, or failing to cover regions that it does.
\end{itemize}

The mathematical objects (loss, gradient, optimizer) are the same. The interpretation is different, and the design choices that follow --- what to model, what to condition on, how to sample --- are different.

\section{What you have seen}

The cross-entropy loss from Day 2 is negative log-likelihood. Maximum likelihood is approximate KL minimization between data and model. The VAE optimizes a tractable lower bound on the log-likelihood and uses a reparameterized sampling trick to backprop through randomness. Diffusion models bypass the likelihood entirely, training a denoiser by plain MSE and sampling by reversing the noising process.

These three ideas --- likelihood, ELBO, denoising score --- cover essentially every generative model used in production: language models (next-token likelihood), Stable Diffusion (denoising), image VAEs (the perceptual compression stages of latent diffusion). Day 4 returns to a question we have ducked: when the data has structure (graphs, sequences, images on a regular grid), what architectures encode that structure into the model?

\section*{Lab}

Open the Day 3 notebook (\texttt{Day3\_Probability\_GenerativeModels.ipynb}). You will:

\begin{enumerate}
  \item Sample from a Gaussian, fit one to data, and compute KL between two Gaussians by hand and by Monte Carlo --- so you see why the closed-form matters.
  \item Train a small VAE on MNIST. Visualize the 2D latent space, sample from it, and decode a path between two latents to watch one digit morph into another.
  \item Implement a denoising diffusion model from scratch on a 2D spiral dataset. Visualize the forward noising process and the reverse generation process step by step.
  \item Train a small diffusion model on MNIST, sample a few digits, and inspect the intermediate noise levels.
\end{enumerate}

By the end of the lab you will have built two generative models from scratch and trained both successfully on real data, using only NumPy, PyTorch, and the math from this chapter.

\chapter{Modelling for AI}
\label{ch:day4}

\begin{aitip}
This chapter is forthcoming. The lecture is outlined below; the full text will be released as part of the next workshop iteration. The syllabus and lab notebook for this day are already available.
\end{aitip}

\section*{Planned sections}
\begin{itemize}
  \item Constrained and unconstrained optimization; regularization as soft constraints.
  \item Invariance and equivariance: what symmetries should the model respect?
  \item Specialized architectures (CNNs, GNNs, attention) as inductive bias.
  \item Translating a real-world problem into a learning problem: choice of inputs, targets, loss, evaluation metric.
\end{itemize}

\section*{Lab}
The Day 4 notebook (\texttt{Day4\_Modeling\_InductiveBias.ipynb}) compares an MLP, a CNN, and a small GNN on a graph-structured task, showing how inductive bias changes sample efficiency.

\chapter{Evaluating AI models}
\label{ch:day5}

\begin{aitip}
This chapter is forthcoming. The lecture is outlined below; the full text will be released as part of the next workshop iteration. The syllabus and lab notebook for this day are already available.
\end{aitip}

\section*{Planned sections}
\begin{itemize}
  \item Basic statistics, the bias--variance decomposition, the meaning of overfitting.
  \item Validation strategies, calibration, uncertainty quantification.
  \item Adversarial examples and out-of-distribution data --- why a lab-perfect model fails in production.
  \item The decision framework: when to ship a model, when to collect more data, when to redesign the problem.
\end{itemize}

\section*{Lab}
The Day 5 notebook (\texttt{Day5\_Evaluation\_Generalization.ipynb}) compares several model architectures under controlled distribution shifts, building intuition for which failures generalize and which do not.


% ============================================================
\backmatter

\chapter*{Appendix A: Environment setup}
\addcontentsline{toc}{chapter}{Appendix A: Environment setup}

\section*{Option 1: Google Colab (recommended)}
Go to \url{https://colab.research.google.com}. Select \textbf{Runtime > Change runtime type > T4 GPU} for the diffusion-model labs in Day 3 and Day 4. All libraries can be installed in notebook cells with \texttt{pip install}.

\section*{Option 2: Local installation}
Python 3.10 or later. For Day 3's diffusion lab and Day 4's graph-learning lab, an NVIDIA GPU with 6+ GB VRAM accelerates training, but every lab can run on CPU in under 15 minutes.

\section*{Core libraries}
\begin{minted}{bash}
pip install torch torchvision numpy matplotlib scikit-learn
pip install jupyter notebook
\end{minted}

\chapter*{Appendix B: Reading list}
\addcontentsline{toc}{chapter}{Appendix B: Reading list}

\begin{itemize}
  \item Goodfellow, Bengio, Courville, \emph{Deep Learning} (MIT Press, free online).
  \item Bishop \& Bishop, \emph{Deep Learning: Foundations and Concepts} (Springer, 2024).
  \item Murphy, \emph{Probabilistic Machine Learning} (MIT Press, free online).
  \item Strang, \emph{Linear Algebra and Learning from Data} (Wellesley-Cambridge, 2019).
  \item Ho, Jain, Abbeel, ``Denoising Diffusion Probabilistic Models'', NeurIPS 2020, \url{https://arxiv.org/abs/2006.11239}.
  \item Vaswani et al., ``Attention Is All You Need'', NeurIPS 2017, \url{https://arxiv.org/abs/1706.03762}.
\end{itemize}

\end{document}
